\section{合成数据与 Mid-training}

\subsection{Mid-training 的定位}

\begin{importantbox}{曹玉：Mid-training 是 World Model 建模的训练中间态}
Mid-training 的定义一直不严谨。早期它是``达不到 SFT Golden Label 质量的指令数据''。现在的认知更高：
\begin{itemize}
  \item RL Post-training 是对已有知识技能的\textbf{组合与泛化}
  \item 这些技能需要在 Mid-training 阶段被模型\textbf{理解和体会}
  \item Mid-training 数据不仅教模型``怎么做事''（Policy Model），也在\textbf{建模世界的状态转移}（World Model）
  \item 观察数学等 Mid-training 数据：Loss Mask 与 SFT 不同，模型可以看到中间状态转移
\end{itemize}
在 Policy Model 和 World Model 功能压缩到同一模型的时代，Mid-training 质量很大程度上决定了 Post-training 质量。
\end{importantbox}

\subsection{合成数据的两大挑战}

\begin{warningbox}{薛福昭：合成数据的核心风险}
\textbf{挑战一：跨代迁移}——当前模型生成的合成数据，其 gain 能否 transfer 到下一代更大的模型？直觉上就是 challenge，实践中也是。相当于用当前模型的 ceiling 去训练下一代。

\textbf{挑战二：前置质量 vs Post-training Gain}——将 Post-training 模型生成的数据``回灌''到更早的训练阶段，可能会蚕食 Post-training 本身的 gain。因为 Pre $\to$ Post 本质是 data quality 不断提升的过程，提前引入高质量数据后，Post-training 的边际收益可能下降。

\textbf{附加风险}：大模型记忆力强，合成数据不可避免包含与训练数据相似的分布，导致隐式 repeat $\to$ overfitting $\to$ Model Collapse。
\end{warningbox}

\subsection{Mid-training 是否引入新知识？}

\textbf{谢天宝}：业界共识是 Post-training 将 Pre/Mid-training 中已存在的 trajectory 的 pass rate 从 pass@64 拉到 pass@1。如果 Pre/Mid-training 没有见过足够好的数据，Post-training 很难强行学到。同时 Post-training 的 token 窗口有限，不能放太多数据，必须在 Mid-training 补回来。

\textbf{刘子伟}：学界有 NeurIPS Best Paper Runner-up 级别的工作研究这个问题——Pre-training 奠定知识的大分布，Mid-training 把 boundary 往外推。但如果考虑原生多模态模型，Mid-training 可能需要不一样的 signal（不只是 language），以及 learning to learn 的能力。

\begin{knowledgebox}{数据分析的科学化}
谢天宝指出：很多厂商和学校对数据的分析能力有限——``数据要不要、放在哪个阶段''主要靠人判断。Active Learning、Self-Improvement 等概念虽然存在，但与工程师手动管理主流模型训练之间仍有很大脱节。数据管理的科学化是一个值得深入研究的方向。
\end{knowledgebox}

\subsection{多模态数据瓶颈与 Vision-Centric 路线}

\textbf{曹玉}：语言模型得益于互联网海量文本，但\textbf{不可能有第二个互联网}提供同等规模的 text-vision pair。低等动物没有 text 概念但生存了很久——VLM 可能不需要那么多数据就能打底，更多应该通过\textbf{环境交互}（RL 或 reward-free 的 learning from experience）来学习。

\textbf{刘子伟}：可以回看前深度学习时代的\textbf{视觉自监督}工作（何恺明、Alusha 等）获得启发。团队尝试过 Visual Jigsaw（打乱图像/视频让模型恢复顺序），能激发 low-level 和 high-level 特性，包括\textbf{空间智能}——这是当前最强模型（GPT-5、Gemini）也做不好的方向。

\begin{importantbox}{薛福昭：Vision-Centric 是长期方向}
反对纯仿生路线（人/动物的 vision perception 很大程度来自 physical touching，机器很难获得这种 signal）。但\textbf{摆脱对文本的依赖}是 make sense 的：pixel reconstruction、next frame prediction 等方式可以提供海量 vision token。当 compute 再翻 10--100 倍时，text-image pairs 一定跟不上，直接用 vision 硬搞是可行的。
\end{importantbox}

\subsection{本章小结}
Mid-training 不只是``质量略低的 SFT''，而是 World Model 建模的训练中间态。合成数据面临跨代迁移和 Model Collapse 风险。多模态数据瓶颈可能通过 Vision-Centric 自监督路线突破。

%% ===== 后训练：RL =====
